\documentclass[DIN, pagenumber=false, fontsize=11pt, parskip=half]{scrartcl}
\usepackage[UTF8]{ctex}  % 中文支持
% \usepackage{ngerman}
\usepackage[utf8]{inputenc}
\usepackage[T1]{fontenc}
\usepackage{textcomp}
\usepackage{graphicx}
\usepackage{amsmath}
\usepackage{booktabs}

\usepackage{xcolor} % Required for custom red

\usepackage{geometry} % Custom margins

\usepackage{tocloft} % Modifying the table of contents

% for matlab code
% bw = blackwhite - optimized for print, otherwise source is colored
\usepackage[framed,numbered,bw]{mcode}
\usepackage[most]{tcolorbox}
% for other code
%\usepackage{listings}

\setlength{\parindent}{0em}

% set section in CM
\setkomafont{section}{\normalfont\bfseries\Large}

\newcommand{\mytitle}[1]{{\noindent\Large\textbf{#1}}}
\newcommand{\mysection}[1]{\textbf{\section*{#1}}}
\newcommand{\mysubsection}[2]{\romannumeral #1) #2}

\newtcolorbox{parchmentquote}{
    enhanced,
    colback=orange!5!white,   % 浅米黄色背景
    colframe=orange!50!black,
    boxrule=0.5pt,
    arc=2pt,                  % 细微圆角
    boxsep=5pt,
    before skip=10pt,
    after skip=10pt,
    fontupper=\itshape,
    borderline west={2pt}{0pt}{orange!80!black} % 加厚的左装饰线
}
%===================================
\begin{document}

\noindent\textbf{课后总结} \hfill \textbf{online course}\\
primary school \hfill Liu\\

\mytitle{长方体与正方体 \hfill \today}
\newline
\begin{parchmentquote}
    熟练掌握长方体、正方体的表面积与体积计算公式。
  正方体是特殊的长方体。

\end{parchmentquote}


%===================================
\section{简单认识长方体与正方体}
\begin{tikzpicture}[scale=1]

% 前面矩形
\draw (0,0) rectangle (4,3);

% 后面矩形（平移）
\draw (1.5,1.2) rectangle (5.5,4.2);

% 连接对应顶点
\draw (0,0)--(1.5,1.2);
\draw (4,0)--(5.5,1.2);
\draw (4,3)--(5.5,4.2);
\draw (0,3)--(1.5,4.2);

% 顶点标记
% \node[below left] at (0,0) {$A$};
% \node[below right] at (4,0) {$B$};
% \node[above right] at (4,3) {$C$};
% \node[above left] at (0,3) {$D$};

% \node[below left] at (1.5,1.2) {$E$};
% \node[below right] at (5.5,1.2) {$F$};
% \node[above right] at (5.5,4.2) {$G$};
% \node[above left] at (1.5,4.2) {$H$};
\node at (2.3,-0.5) {长方体};

\end{tikzpicture}
\quad \quad
\begin{tikzpicture}[scale=1.05]

% 前面正方形
\draw (0,0) rectangle (3,3);

% 后面正方形（平移）
\draw (1.2,1.2) rectangle (4.2,4.2);

% 连接对应顶点
\draw (0,0)--(1.2,1.2);
\draw (3,0)--(4.2,1.2);
\draw (3,3)--(4.2,4.2);
\draw (0,3)--(1.2,4.2);

% 顶点标记
% \node[below left] at (0,0) {$A$};
% \node[below right] at (3,0) {$B$};
% \node[above right] at (3,3) {$C$};
% \node[above left] at (0,3) {$D$};

% \node[left] at (1.2,1.2) {$E$};
% \node[right] at (4.2,1.2) {$F$};
% \node[above right] at (4.2,4.2) {$G$};
% \node[above left] at (1.2,4.2) {$H$};
\node at (2.1,-0.5) {正方体};
\end{tikzpicture}
\section{长方体与正方体的性质}

\subsection*{一、长方体的性质}

设长方体的长、宽、高分别为 $a,b,c$。

\begin{enumerate}
    \item 有 $6$ 个面，且每个面都是长方形；
    \item 相对的面互相平行且全等；
    \item 有 $12$ 条棱；
    \item 相对的棱平行且相等；
    \item 有 $8$ 个顶点；
    \item 每个顶点处有三条棱互相垂直；
    \item 体对角线相等，并且交于同一点；
    \item 体积公式：
    \[
    V=abc
    \]
    \item 表面积公式：
    \[
    S=2(ab+bc+ac)
    \]
    % \item 体对角线长度：
    % \[
    % d=\sqrt{a^2+b^2+c^2}
    % \]
\end{enumerate}

\bigskip

\subsection*{二、正方体的性质}

设正方体棱长为 $a$。

\begin{enumerate}
    \item 有 $6$ 个完全相同的正方形面；
    \item 有 $12$ 条棱，且所有棱长都相等；
    \item 有 $8$ 个顶点；
    \item 每个面都是正方形；
    \item 每个顶点处三条棱两两垂直；
    \item 所有面对角线相等；
    \item 所有体对角线相等；
    \item 体积公式：
    \[
    V=a^3
    \]
    \item 表面积公式：
    \[
    S=6a^2
    \]
    % \item 面对角线长度：
    % \[
    % d_{\text{面}}=a\sqrt{2}
    % \]
    % \item 体对角线长度：
    % \[
    % d_{\text{体}}=a\sqrt{3}
    % \]
\end{enumerate}

\bigskip
\section{练习题}
\begin{enumerate}
    \item 已知一个长方体的长为 $5$ cm，宽为 $3$ cm，高为 $2$ cm，求它的体积和表面积。
    \item 一个正方体的棱长为 $0.15$ cm，求它的体积和表面积。
    \item 如果一个长方体的体积是 $60 cm^3$，且长、宽分别是 $5$ cm和 $4$ cm，那么这个长方体的表面积是多少？
    \item 已知一个正方体的表面积是 $9.6 cm^2$，求它的棱长和体积。
    \item 一个长方体的长是$2.4cm$,宽是长的$\frac{1}{2}$，高是宽的$\frac{3}{4}$，求它的体积。
\end{enumerate}     

\end{document}